\documentclass[10pt, twocolumn, letterpaper]{article}

\usepackage{times}
\usepackage{graphicx}
\usepackage{amsmath}
\usepackage{amssymb}
\usepackage{booktabs}
\usepackage{hyperref}
\usepackage{geometry}
\geometry{margin=1in}

\title{}
\author{Vayun Gupta \\ University of Illinois Urbana-Champaign}

\usepackage{xcolor}

\newcommand{\ragini}[1]{\textcolor{red}{\textbf{[Ragini: #1]}}}
\begin{document}

\maketitle

\begin{abstract}

\end{abstract}

\section{Introduction}
Key points
IoT systems rely on location information for critical services
Sensors can misreport (spoof) their location
Existing approaches:
GPS (not always reliable / spoofable)
RF-based localization (noisy, environment-sensitive)
Challenge:
RF signals are affected by multipath, obstruction, and noise
Hard to distinguish:
true environmental distortion vs malicious spoofing

\section{Motivations}

Your idea-
Use multi-modal verification:
RF (RSSI, SNR, trilateration)
Environmental context (satellite → buildings, trees)
Temporal consistency
Move from:
single-shot localization → trust verification problem


\section{Background}

\subsection{RF based Localization}
1. Trilateration using RSSI/SNR
2. Distance estimation from signal strength
Limitations:
multipath fading
NLOS conditions
AP geometry sensitivity

\subsection{Environmental Impact on RF Signals}
Talk about SatLoc paper, Line of sight vs. non line of sight conditions, motivation for using satellite images


\subsection{Threat Model
}
Define clearly (very important):

Sensor may:
falsify location
manipulate signal properties (RSSI/SNR)
\textbf{Adversary goal:}
appear at a false location while remaining plausible


Assumptions:
1. APs are trusted
2. satellite imagery is reliable attacker cannot perfectly emulate all modalities

\subsection{Using Multi agent system for sensor location verification}
** Need to talk about the different papers that are using agents in sensors verification, localization and cybersecurity 




\section{System Design}
Have a figure here from the Virginia dataset and what the different modules that run on sensor and access points, the role of sensor, access points and agents (what different agents are used). Take reference from safecoop paper for the system design. 

Our system follows a two-stage verification design. First, deterministic RF modules compute packet-count anomalies, RSSI shifts, and trilateration residual changes from gateway observations. These modules remain the calibration anchor because they are reproducible, directly tied to the measurement process, and easy to validate. Second, an LLM-based supervisor is invoked only when those numeric signals are ambiguous, mutually inconsistent, or close to a decision boundary. In this role, the LLM does not replace localization or anomaly scoring; it adjudicates between competing explanations and decides whether the current evidence is sufficient for a trust decision.

This distinction is important for the overall contribution. If a simple threshold already separates attack from benign behavior, the heuristic detector is preferable. The LLM becomes justified only when trust verification requires multi-step reasoning over partially conflicting evidence rather than a single clean scalar threshold.


Describe architecture:

** Supervisor agent
** RF agent
** Environment agent
** Temporal agent
** (optional) verification strategy agent

Pipeline:

Claim → Agents → Tool Calls → Evidence → Iterative Reasoning → Trust Decision

\subsection{Baseline: RF-based Localization}
3.1 Trilateration Model
Convert RSSI → distance
Use multiple APs for triangulation
Solve for location estimate

You can include:

least squares formulation. 


Key argument for agents later:

** RF mismatch is not equal to spoofing
** environmental effects (NLOS) cause false alarms
cannot reason about:
** why mismatch occurs
** whether mismatch is explainable
\subsection{Agentic Trust Verification Framework}

** This is your main novelty section.

--> Motivation for Agents
Localization is not a one-shot problem
Requires:
1. iterative reasoning
2. conflict resolution
3. adaptive evidence collection

Connect to idea:

similar to tool-using agents in complex environments

In our implementation, the LLM is best described as an ambiguity resolver rather than a heuristic replacement. The deterministic layer first summarizes the scenario into compact structured evidence such as top sensor anomaly score, top gateway anomaly score, packet-retention ratios, and trilateration error deltas. The LLM receives only this structured summary and is constrained to produce a bounded output: final label, rationale, cited evidence keys, and whether more evidence should be requested. This keeps the numeric tools as the source of truth and prevents the LLM from inventing measurements or bypassing the underlying RF pipeline.


\subsection{Tool Abstraction}

Define tools clearly:

** RF tools:
1. trilateration
2. AP subset selection


** Environment tools:
1. building/tree extraction
2. LOS estimation



** Verification tools:
1. request additional AP/time window

This is where you justify tool calling.

\subsubsection{Agent Roles}
** RF Agent
1. analyzes signal evidence
2. computes mismatch and confidence


** Environment Agent
1. explains RF behavior via obstacles
2. challenges RF conclusions


** (optional) Temporal Agent
1. checks trajectory plausibility


** Supervisor Agent
1. coordinates agents
2. resolves conflicts
3. decides next action


** Agent Interaction Protocol

1. Important for systems paper:

a) agents exchange:
b) confidence scores
c) explanations


** supervisor:
1. selects next agent/tool
decides when to stop

\subsection{Sequential Verification Process}

Step-by-step:

Evaluate RF evidence
If mismatch → check environment\\

If still uncertain → check temporal consistency\\

If ambiguity remains → gather more data\\

Final trust decision\\


\subsection{Decision Policy}
1. Accept

2. Reject

3. Suspicious

4. Request more evidence

Operationally, the heuristic stage emits one of three outcomes: accept, reject, or uncertain. The LLM is triggered only for uncertain cases or when the evidence suggests multiple plausible explanations, such as overlapping sensor-side and gateway-side anomalies. If the LLM is unavailable, malformed, or low-confidence, the system deterministically falls back to the heuristic decision.

\section{Experimental Setup}

\subsubsection{Datasets}
Talk about following:
** AP locations
** sensor data RSSI/SNR
** satellite imagery
** synthetic spoofing attacks

\subsection{Baselines}
One baseline is trilateration, second baseline is having trust verification system without any agents at all.

To isolate the value of the LLM, we evaluate three progressively richer settings: pure heuristic detection, heuristic detection with explanation-only LLM support, and heuristic detection with bounded LLM adjudication. This avoids framing the problem as ``LLM versus heuristic'' and instead tests whether the LLM adds value specifically on hard trust-verification cases.

\subsubsection{Attack Scenarios}
** random location spoofing
** environment aware spoofing (?) ignore if we cannot or put in future works

\subsubsection{Metrics}

** Metrics
** detection accuracy
** false positive rate
** localization error *you already have the cdf plots)
** trust calibration
** number of extra verification steps (agent efficiency)

\section{Results and Analysis}

1. RF based localization \\
2. Impact of Environment Features
improvement over RF \\
3. Agentic Framework Performance \\ accuracy improvements
better handling of ambiguous cases

4. Efficiency Analysis
fewer unnecessary verifications
cost vs accuracy tradeoff plot with agents (can you check the related papers what other metrics they have used ????)

The most important slice is not overall accuracy alone, but performance on ambiguous cases where deterministic cues conflict or lie near threshold values. In addition to standard detection accuracy, we therefore report ambiguous-case accuracy, false positive rate on clean and weak-noise benign scenarios, LLM invocation rate, and explanation-faithfulness rate. These metrics show whether the LLM is selective rather than always-on and whether its benefits come from resolving genuinely hard cases rather than replacing well-calibrated numeric detectors.


\subsection{Case Studies}

show examples where:
** RF fails and agents resolve ambiguity





\subsection{Discussion}

The main argument is not that the LLM is universally better than heuristics. Rather, heuristics are better for stable low-level measurement and should remain the primary detector. The LLM becomes useful only when the system must reason across heterogeneous evidence, explain why RF mismatch may or may not be malicious, or recommend gathering more evidence before making a trust decision.


\subsubsection{When Agents Help Most}
-> NLOS environments
-> sparse AP settings
-> ambiguous RF signals


\subsubsection{Failure Modes}


--> insufficient exploration
--> wrong tool choice
--> over-reliance on one modality


\section{Limitations}
--> dependence on satellite accuracy
--> latency due to sequential reasoning
\section{Future Work}




\section*{Acknowledgment}

\section*{References}

\end{document}
